\color{unimain}{\subsection{Автоматика опережающего деления сети}}
\color{black}

\begin{enumerate}

\item
В качестве мер по ограничению токов КЗ применятся автоматика опережающего деления сети (АОДС).
При использовании АОДС в интервале между возникновением КЗ и отключением поврежденного
присоединения отключается ШСВ, а после отключения поврежденного присоединения при необходимости
ШСВ автоматически включаются в цикле АПВ. Это позволяет уменьшить ток КЗ, отключаемый выключателем
поврежденного присоединения. Указанная операция осуществляется, как правило, без задержки отключения
выключателя поврежденного присоединения благодаря использовании быстродействующих измерительных
органов АОДС и выключателя с малым временем отключения.

\item
Токовые измерительные органы АОДС основаны на принципе замера периодической составляющей тока по
мгновенным значениям его первой производной в интервале точной трансформации ТТ. Совпадение знаков
тока и его первой производной является устойчивым признаком, позволяющим отстроиться от искажений тока
при насыщения магнитопровода ТТ.

\item
Функциональная схема АОДС представлена на рисунке \ref{pict_aods:AODS}.

\begin{figure}[H]
\centering
\includegraphics[width=0.8\textwidth,height=0.8\textheight,keepaspectratio]{\fbpath/08.01 СВ/661.3051 АОДС/img/AODS.pdf}
\captionof{figure}{Функциональная схема АОДС}\label{pict_aods:AODS}
\end{figure}

\item
Порог тока срабатывания АОДС, реагирующий на фазный ток, задается уставкой \textbf{«Iф АОДС»}. Действие
ступеней осуществляется без выдержки времени на отключение ШСВ. Ввод в работу осуществляется с
помощью программного ключа \textbf{«АОДС»}. 

\item
Оперативный ввод/вывод АОДС производится с помощью БУ (блоков управления функциями) \textit{«АОДС СВ/АОДС: Вывод»}.

\item
Уставочные параметры функции АОДС представлены в таблице \ref{app_set:aods}.

\end{enumerate} 